% Generated draft: requires mathematical and source review.
\documentclass[11pt]{article}
\usepackage[margin=25mm]{geometry}
\usepackage{fontspec}
\setmainfont{Latin Modern Roman}
\usepackage{amsmath,amssymb,mathtools,unicode-math}
\setmathfont{Latin Modern Math}
\usepackage{enumitem,needspace,xurl,hyperref}
\setlength{\parindent}{0pt}
\setlength{\parskip}{5pt}
\setlength{\emergencystretch}{3em}
\newcommand{\N}{\mathbb{N}}
\newcommand{\Z}{\mathbb{Z}}
\newcommand{\Q}{\mathbb{Q}}
\newcommand{\R}{\mathbb{R}}
\newcommand{\C}{\mathbb{C}}
\newcommand{\F}{\mathbb{F}}
\newcommand{\A}{\mathbb{A}}
\newcommand{\PP}{\mathbb{P}}
\newcommand{\E}{\mathbb{E}}
\newcommand{\eps}{\varepsilon}
\newcommand{\dd}{\,\mathrm{d}}
\begin{document}
\section*{UnsolvedMath ID: 30000459}
\textbf{Source problem code:} catalogue entry
\subsection{Definitions and mathematical statement}
Let $B$ be a block of a symmetric-group algebra $k\mathfrak S_n$ in the report's modular setting, with abelian defect group. Let $D_\lambda,D_\mu$ be simple modules in that same block, indexed by its admissible partition convention, and let $\delta_e$ be the partition invariant used in the source. The conjectural implication is
\[
 \delta_e(\lambda)=\delta_e(\mu)
 \quad\Longrightarrow\quad
 \operatorname{Ext}^1_B(D_\lambda,D_\mu)=0.
\]
Equivalently, every short exact sequence $0\to D_\mu\to E\to D_\lambda\to0$ splits. Ext may equally be computed in the ambient group-algebra module category because both modules lie in the same block.

The extract does not define $e$ or $\delta_e$, fix characteristic or splitting assumptions on $k$, or specify restricted versus regular partition labels. No guessed parity formula or equality of cores is substituted for the exact invariant. It also does not require $\lambda\ne\mu$, so the literal formula includes self-extensions.

The characteristic qualification is substantive. If characteristic two and equal labels are included, the unique simple module of $k\mathfrak S_2\cong k[t]/(t^2)$ has nonzero self-$\operatorname{Ext}^1$, while its block has abelian defect group and its invariant equals itself. Thus the universal wording must be matched to the intended source restrictions, for example any odd-characteristic or label qualifications actually imposed there, rather than declared valid in all characteristics.

The literature assessment checked on 21 August 2026 reports verification for all weight-three blocks in the intended setting, while saying no proof of the general assertion was found. Those reported cases must retain their original characteristic/block hypotheses; they do not remove the missing scope or define $\delta_e$. The partially-solved label is dated supplied triage requiring independent verification.
\subsection{Short English statement}
In an abelian-defect symmetric-group block, does equality of the source's partition invariant force first Ext to vanish? Its invariant and characteristic restrictions are missing; the literal all-characteristic version includes a characteristic-two self-extension obstruction.
\subsection{Sources}
\begin{itemize}
\item \emph{Representations of Finite Groups} (2007), Oberwolfach Reports, DOI 10.4171/owr/2006/15, pp. 895--978, \url{https://doi.org/10.4171/owr/2006/15}; proposed year supplied as 2006.
\item The sole source in the literature assessment checked 21 August 2026 is the original report, \url{https://doi.org/10.4171/owr/2006/15}, which is reported to verify the weight-three cases.
\item Local UnsolvedMath registry, permanent ID $30000459$, exact source code \texttt{catalogue entry}, created 25 August 2026 and updated 16 September 2026. Supplied verification preserves the displayed implication while normalizing module, partition-invariant and Ext notation.
\end{itemize}
\end{document}
